\documentclass[11pt]{amsart}

\usepackage[T1]{fontenc}
\usepackage{lmodern}
\usepackage{microtype}
\usepackage{mathtools,amssymb,amsthm}
\usepackage{array}
\usepackage[hidelinks]{hyperref}

\hypersetup{
  pdftitle={Contractibility Can Fail at the Top of the Rank-Spread Filtration:
            A Ten-Element Graded Poset}
}

\newtheorem{theorem}{Theorem}
\newtheorem{lemma}[theorem]{Lemma}
\newtheorem{proposition}[theorem]{Proposition}

\DeclareMathOperator{\St}{St}
\newcommand{\rk}{\mathrm{rk}}
\newcommand{\Dl}{\Delta}
\newcommand{\Susp}{\Sigma}
\newcommand{\ZZ}{\mathbb Z}
\newcommand{\Ho}{\widetilde H}

\title[Contractibility can fail at the top filtration stage]
{Contractibility Can Fail at the Top of the Rank-Spread Filtration:
A Ten-Element Graded Poset}

\date{}

\subjclass[2020]{06A07, 55U10, 05E45}
\keywords{graded poset, order complex, rank-spread filtration, contractibility,
collapsibility, magnitude homology, counterexample}

\begin{document}

\begin{abstract}
Let $P$ be a finite graded poset of rank $n$ and let $\Dl^{(k)}(P)$ be the
subcomplex of the order complex $\Dl(P)$ spanned by the chains of rank spread
at most $k$, so that $\Dl^{(0)}(P)$ is the vertex set and
$\Dl^{(n)}(P)=\Dl(P)$. Kitajima asks whether $\Dl^{(k)}(P)$ is contractible
whenever $\Dl^{(k-1)}(P)$ is. We exhibit a counterexample to that implication
as stated. Let $P$ be the poset on the ten elements
$a_0,a_1,b_0,b_1,b_2,c_0,c_1,c_2,d_0,d_1$, of ranks $0,0,1,1,1,2,2,2,3,3$,
whose covers are exactly the fifteen relations displayed in
Section~\ref{sec:witness}. Then $P$ is graded of rank $3$, the complex
$\Dl^{(2)}(P)$ collapses to a point in $27$ elementary steps, and
$\Dl^{(3)}(P)=\Dl(P)$ is homotopy equivalent to $S^2$; in particular
$\Ho_2(\Dl^{(3)}(P);\ZZ)=\ZZ$ and $\Dl^{(3)}(P)$ is not contractible. The
witness has $k=n$, so what fails to be contractible is the full order complex;
whether the implication can fail at a proper stage $k<n$ we leave open.
\end{abstract}

\maketitle

\section{The question}\label{sec:question}

Let $P$ be a finite poset. Its \emph{order complex} $\Dl(P)$ is the simplicial
complex whose faces are the nonempty chains of $P$. Following
Kitajima~\cite{Kitajima}, $P$ is \emph{graded of rank $n$} if every maximal
chain of $P$ has length $n$ --- that is, $P$ is \emph{pure}; no bound is
imposed, so $P$ may have several minimal and several maximal elements. A graded
poset carries a rank function $\rk\colon P\to\{0,\ldots,n\}$, the
\emph{rank spread} of a chain $C\subseteq P$ is
$\max_{x\in C}\rk(x)-\min_{x\in C}\rk(x)$, and for $0\le k\le n$
\[
  \Dl^{(k)}(P)\ :=\ \bigl\{\,C\in\Dl(P)\ :\ \text{the rank spread of $C$ is at
  most }k\,\bigr\}
\]
is a subcomplex of $\Dl(P)$. Thus $\Dl^{(0)}(P)$ is the vertex set,
$\Dl^{(1)}(P)$ is the Hasse diagram viewed as a $1$-complex,
\[
  \Dl^{(0)}(P)\subseteq\cdots\subseteq\Dl^{(n)}(P)=\Dl(P),
\]
and the filtration is the poset case of the filtered nerve of
Di et al.~\cite{DIMZ}, whose relative homology
$H_*\bigl(\Dl^{(k)}(P),\allowbreak\Dl^{(k-1)}(P)\bigr)$ is the magnitude
homology of $P$
\cite[Lemma~1.7]{DIMZ}; the order-complex antecedent for metric spaces is due
to Kaneta and Yoshinaga~\cite{KY}. In that context Kitajima records, verbatim:

\begin{quote}
``It is an open problem whether $\Dl^{(k)}(P)$ is contractible whenever
$\Dl^{(k-1)}(P)$ is contractible.''
\end{quote}

The sentence is in Remark~2.13 on page~8 of \cite{Kitajima} (version~v2 of
24 June 2026); the ambient definition of ``graded of rank $n$'' as purity is
Definition~1.9 on page~5.

Since the sentence quantifies over all finite graded posets $P$, all $n$, and
all $k$, a single pair $(P,k)$ with $\Dl^{(k-1)}(P)$ contractible and
$\Dl^{(k)}(P)$ not contractible settles it, negatively. The witness below has
$k=n$; Section~\ref{sec:scope} says what that does and does not leave open.

\section{The witness}\label{sec:witness}

Let $P$ be the poset on the ten elements
\[
  \underbrace{a_0,\ a_1}_{\rk=0}\qquad
  \underbrace{b_0,\ b_1,\ b_2}_{\rk=1}\qquad
  \underbrace{c_0,\ c_1,\ c_2}_{\rk=2}\qquad
  \underbrace{d_0,\ d_1}_{\rk=3}
\]
whose cover relations are exactly the following fifteen:
\begin{equation}\label{eq:covers}
\begin{array}{lll}
 a_0<b_0 & &\\
 a_1<b_1 & a_1<b_2 &\\
 b_0<c_0 & b_0<c_1 & b_0<c_2\\
 b_1<c_0 & b_1<c_1 &\\
 b_2<c_0 & b_2<c_2 &\\
 c_0<d_0 & &\\
 c_1<d_0 & c_1<d_1 &\\
 c_2<d_0 & c_2<d_1 &
\end{array}
\end{equation}
The order on $P$ is the transitive closure of \eqref{eq:covers}. It consists of
$31$ strict pairs: the fifteen covers, the six pairs $a_i<c_j$ that are
present, the six pairs $b_i<d_j$, and all four pairs $a_i<d_j$. Recomputing
the cover relation from that closure returns exactly the fifteen pairs of
\eqref{eq:covers}, so \eqref{eq:covers} is a faithful presentation.

\begin{lemma}\label{lem:graded}
$P$ is graded of rank $3$; it has two minimal elements $a_0,a_1$ and two
maximal elements $d_0,d_1$, hence neither a minimum nor a maximum, and
$\Dl^{(3)}(P)=\Dl(P)$.
\end{lemma}

\begin{proof}
Every relation in \eqref{eq:covers} raises the rank by exactly $1$; every
element of rank $<3$ has an upper cover and every element of rank $>0$ has a
lower cover. Hence every maximal chain of $P$ starts at rank $0$, ends at
rank $3$ and passes through one element of each rank, so it has length $3$;
there are eleven such chains, and $P$ is graded of rank $3$. Only $a_0,a_1$
have no lower cover and only $d_0,d_1$ have no upper cover, so these are
exactly the minimal and the maximal elements; each of the two sets has two
members, so $P$ has neither a minimum nor a maximum. Finally, every chain of
$P$ has rank spread at most $3$, so $\Dl^{(3)}(P)=\Dl(P)$.
\end{proof}

The face numbers of the two filtration stages used below are
\begin{equation}\label{eq:fvec}
  \Dl^{(2)}(P)\colon\ (10,27,18),\qquad
  \Dl^{(3)}(P)\colon\ (10,31,34,11),
\end{equation}
with Euler characteristics $+1$ and $+2$.

For $S\subseteq\{0,1,2,3\}$ write $P_S=\{x\in P:\rk(x)\in S\}$ and abbreviate
$P_{[i,j]}=P_{\{i,i+1,\ldots,j\}}$.

We use one standard fact throughout.

\begin{lemma}\label{lem:susp}
Let $X=A\cup B$ be a simplicial complex covered by two subcomplexes with
$Y=A\cap B\ne\emptyset$. If $A$ and $B$ are contractible, then $X$ is homotopy
equivalent to the suspension $\Susp Y$.
\end{lemma}

\begin{proof}
$X$ is the pushout of the two inclusions $A\hookleftarrow Y\hookrightarrow B$,
which are cofibrations of CW pairs, so $X$ is also their homotopy pushout.
Replacing $A$ and $B$ by points, which is a homotopy equivalence of pushout
diagrams, gives $X\simeq\bigl(\{*\}\leftarrow Y\rightarrow\{*\}\bigr)$-pushout
$=\Susp Y$.
\end{proof}

For a vertex $v$ of a complex $X$ we write $\St_X(v)$ for its closed star, the
subcomplex of faces $\sigma$ with $\sigma\cup\{v\}\in X$, together with their
faces. If $X=\Dl(Q)$ is an order complex and $v\in Q$, then $\St_X(v)$ is the
cone with apex $v$ over the order complex of $\{x\in Q:x<v\text{ or }x>v\}$,
so it is contractible.

\section{$\Dl^{(3)}(P)$ is homotopy equivalent to $S^2$}\label{sec:top}

\begin{proposition}\label{prop:S2}
$\Dl^{(3)}(P)=\Dl(P)\simeq S^2$. In particular
$\Ho_2(\Dl(P);\ZZ)=\ZZ$ and $\Dl(P)$ is not contractible.
\end{proposition}

\begin{proof}
The maximal elements of $P$ are $d_0$ and $d_1$, and they are incomparable, so
no chain contains both, while every chain of $P$ that is maximal ends in $d_0$
or in $d_1$. Hence
\[
  \Dl(P)=\St(d_0)\cup\St(d_1),
\]
and both stars are cones, hence contractible. Their intersection is the order
complex of $\{x\in P:x<d_0\}\cap\{x\in P:x<d_1\}$. Now
$\{x<d_0\}=P_{[0,2]}$, all eight elements of rank $\le2$, whereas
$\{x<d_1\}=P_{[0,2]}\setminus\{c_0\}$: indeed $c_1,c_2<d_1$ by
\eqref{eq:covers}, hence $b_0,b_1<c_1<d_1$ and $b_2<c_2<d_1$ and
$a_0<b_0<d_1$, $a_1<b_1<d_1$, while $c_0<d_1$ is not among the relations of
$P$. So
\[
  Y:=\St(d_0)\cap\St(d_1)=\Dl\bigl(P_{[0,2]}\setminus\{c_0\}\bigr),
\]
the complex on the seven vertices $a_0,a_1,b_0,b_1,b_2,c_1,c_2$ with the eleven
edges
\[
  a_0b_0,\ a_1b_1,\ a_1b_2,\ b_0c_1,\ b_0c_2,\ b_1c_1,\ b_2c_2,\
  a_0c_1,\ a_0c_2,\ a_1c_1,\ a_1c_2
\]
and the four triangles $a_0b_0c_1$, $a_0b_0c_2$, $a_1b_1c_1$, $a_1b_2c_2$; its
Euler characteristic is $7-11+4=0$.

Apply the same argument to $Y$, at the bottom instead of the top. Every vertex
of $Y$ other than $a_0,a_1$ lies above $a_0$ or above $a_1$, so
$Y=\St_Y(a_0)\cup\St_Y(a_1)$ with both stars cones. In $Y$ we have
$\{x>a_0\}=\{b_0,c_1,c_2\}$ and $\{x>a_1\}=\{b_1,b_2,c_1,c_2\}$, whose
intersection is $\{c_1,c_2\}$, two incomparable elements. So
$\St_Y(a_0)\cap\St_Y(a_1)$ is a pair of points, i.e.\ $S^0$, and
Lemma~\ref{lem:susp} gives $Y\simeq\Susp S^0=S^1$. A second application of
Lemma~\ref{lem:susp}, now to $\Dl(P)=\St(d_0)\cup\St(d_1)$, gives
\[
  \Dl(P)\ \simeq\ \Susp Y\ \simeq\ \Susp S^1\ =\ S^2 .\qedhere
\]
\end{proof}

As a check, $\chi(\Dl(P))=10-31+34-11=2=\chi(S^2)$ by \eqref{eq:fvec},
whereas a contractible complex has $\chi=1$.

\section{$\Dl^{(2)}(P)$ is collapsible}\label{sec:collapse}

Recall that a face $\tau$ of a simplicial complex $X$ is \emph{free} if it is
contained in exactly one proper coface $\sigma$; deleting the pair
$\{\tau,\sigma\}$ is an \emph{elementary collapse} and does not change the
homotopy type. A complex that can be reduced to a single vertex by elementary
collapses is \emph{collapsible}, hence contractible.

\begin{proposition}\label{prop:coll}
$\Dl^{(2)}(P)$ is collapsible, hence contractible.
\end{proposition}

\begin{proof}
$\Dl^{(2)}(P)$ has $10+27+18=55$ faces by \eqref{eq:fvec}. The following
$27$ elementary collapses delete $54$ of them and leave the single vertex
$d_1$. In each row, $\tau$ has exactly one proper coface in the complex
remaining at that step, namely $\sigma$; the reader may check the rows in
order, and each check involves listing the at most three cofaces of a face.
\begin{center}
\renewcommand{\arraystretch}{1.05}
\begin{tabular}{r>{$}l<{$}>{$}l<{$}@{\quad}r>{$}l<{$}>{$}l<{$}@{\quad}r>{$}l<{$}>{$}l<{$}}
 & \tau & \sigma & & \tau & \sigma & & \tau & \sigma\\
\hline
 1 & a_0c_0 & a_0b_0c_0 & 10 & b_0c_0 & b_0c_0d_0 & 19 & b_2 & b_2d_1\\
 2 & a_0c_1 & a_0b_0c_1 & 11 & b_1c_0 & b_1c_0d_0 & 20 & c_1d_0 & b_0c_1d_0\\
 3 & a_0b_0 & a_0b_0c_2 & 12 & b_1d_0 & b_1c_1d_0 & 21 & b_0c_1 & b_0c_1d_1\\
 4 & a_0    & a_0c_2    & 13 & b_1c_1 & b_1c_1d_1 & 22 & c_1 & c_1d_1\\
 5 & a_1c_1 & a_1b_1c_1 & 14 & b_1    & b_1d_1    & 23 & b_0d_0 & b_0c_2d_0\\
 6 & a_1b_1 & a_1b_1c_0 & 15 & b_2c_0 & b_2c_0d_0 & 24 & d_0 & c_2d_0\\
 7 & a_1c_0 & a_1b_2c_0 & 16 & c_0    & c_0d_0    & 25 & b_0c_2 & b_0c_2d_1\\
 8 & a_1b_2 & a_1b_2c_2 & 17 & b_2d_0 & b_2c_2d_0 & 26 & b_0 & b_0d_1\\
 9 & a_1    & a_1c_2    & 18 & b_2c_2 & b_2c_2d_1 & 27 & c_2 & c_2d_1
\end{tabular}
\end{center}
The face counts run $55,53,51,\ldots,3,1$, and the last surviving face is the
vertex $d_1$.
\end{proof}

\section{The result, and what it does not settle}\label{sec:scope}

\begin{theorem}\label{thm:main}
Let $P$ be the ten-element poset of \eqref{eq:covers}. Then $P$ is graded of
rank $3$, $\Dl^{(2)}(P)$ is collapsible --- in particular contractible --- and
$\Dl^{(3)}(P)$ is homotopy equivalent to $S^2$, so it is not contractible.
Consequently it is false that $\Dl^{(k)}(P)$ is contractible whenever
$\Dl^{(k-1)}(P)$ is.
\end{theorem}

\begin{proof}
Lemma~\ref{lem:graded}, Proposition~\ref{prop:coll} and
Proposition~\ref{prop:S2}.
\end{proof}

What falls is the universally quantified implication of Remark~2.13, p.~8 of
\cite{Kitajima}, and nothing more. The witness has $k=n=3$, so the complex that
fails to be contractible is the \emph{full} order complex
$\Dl^{(n)}(P)=\Dl(P)$; that value of $k$ lies in the range $0\le k\le n$ of the
definition recalled in Section~\ref{sec:question} and inside the sentence's
quantifier, so the implication is false as stated. It is not a witness at a
\emph{proper} filtration stage, and this note leaves untouched the question
whether contractibility passes from $\Dl^{(k-1)}(P)$ to $\Dl^{(k)}(P)$ when
$k<n$.

No minimality is claimed: neither that ten is the least number of elements of
such a witness, nor that $k=3$ or $n=3$ is least.

\section{Verification}\label{sec:verify}

Everything asserted above is determined by the fifteen relations
\eqref{eq:covers}, and Sections~\ref{sec:top}--\ref{sec:collapse} are proofs a
reader can follow with no machine. The program \texttt{verify.\allowbreak py}
accompanying this note, with its recorded output in
\texttt{verify.\allowbreak output.\allowbreak txt}, reads
\eqref{eq:covers} as printed and re-derives, in exact integer arithmetic and
using only the Python standard library, the transitive closure, the cover
relation, purity, the $f$-vectors and Euler characteristics of
\eqref{eq:fvec}, the integral homology of the filtration stages by Smith
normal form, and the intermediate objects of the proof of
Proposition~\ref{prop:S2}. The collapse of Section~\ref{sec:collapse} is
checked rather than produced: the $27$ rows displayed in
Proposition~\ref{prop:coll} are transcribed into the program and replayed,
verifying at each row that the printed $\tau$ has exactly one proper coface in
the complex then remaining and that this coface is the printed $\sigma$, and
that after row $27$ the vertex $d_1$ alone remains.

The recorded output reports that all $92$ of its checks pass. It was produced
for a longer draft: its headings follow that draft's numbering, and it also
lists checks and closing notes --- among them a Mayer--Vietoris computation of
the homology of $\Dl^{(2)}(P)$, a discussion of coning this $P$, a census, and
a correction to a remark of \cite{Kitajima} --- that concern material absent
from the present note and on which nothing above depends. In particular
nothing in it bears on the open case $k<n$.

\begin{thebibliography}{9}

\bibitem{DIMZ}
Y.~Di, S.~Ivanov, L.~Mukoseev, and M.~Zhang,
\emph{On the path homology of Cayley digraphs and covering digraphs},
J. Algebra \textbf{653} (2024), 156--199;
\href{https://arxiv.org/abs/2305.15683}{arXiv:2305.15683}.

\bibitem{KY}
R.~Kaneta and M.~Yoshinaga,
\emph{Magnitude homology of metric spaces and order complexes},
Bull. London Math. Soc. \textbf{53} (2021), no.~3, 893--905;
\href{https://arxiv.org/abs/1803.04247}{arXiv:1803.04247}.

\bibitem{Kitajima}
Y.~Kitajima,
\emph{Filtered order complexes and magnitude homology of finite graded
posets},
preprint, 2026.
\href{https://arxiv.org/abs/2606.15241}{arXiv:2606.15241}.
Statement numbers are those printed in version~v2 (24 June 2026), whose
sectioning begins at Section~0.

\end{thebibliography}

\end{document}
